\uuid{JisB}
\titre{Équation droite/plan}
\auteur{bodin ; hanani ; mzari}
\organisation{exo7}
\licence{CC BY-NC-SA 4.0}
\datecreate{2020-01-20}
\niveau{L1}
\module{Géométrie}
\chapitre{Géométrie de l'espace}
\sousChapitre{Plans}
\difficulte{3}
\theme{Plans -- Droites, équation droite, plan}
\exo7qcmid{lille-1.geometrie-espace.020}

\contenu{
\texte{
Pour ces questions, l'espace est muni d'un repère orthonormé direct  $(O, \vec{i}, \vec{j}, \vec{k})$.
}

\begin{question}
\qclassification{141.03,141.04}
\qtags{motcle=équation droite/plan}
Soit  $D_1 : \left\{\begin{array}{ccl}y&=&x+2\\z&=&x \end{array}\right.$,
$D_2 : \left\{\begin{array}{ccl}y&=&2x+1\\z&=&2x-1 \end{array}\right.$ et $\Delta$ une droite parallèle au plan $(xOy)$ et rencontrant les droites $D_1$, $D_2$ et l'axe $(Oz)$.
Quelles sont les assertions vraies ?
\begin{answers}

       \bad {Une équation cartésienne de $\Delta$ est :
     $ \left\{\begin{array}{ccl}y&=&1\\z&=&-1 \end{array}\right.$ ou
     $ \left\{\begin{array}{ccl}x+y+z&=&0\\z&=&1 \end{array}\right.$.}

      \bad {$\Delta$ est contenu dans le plan $z=-1$ ou $z=1$.}

     \good {Une équation cartésienne de $\Delta$ est :
     $ \left\{\begin{array}{ccl}y&=&0\\z&=&-2 \end{array}\right.$ ou
     $ \left\{\begin{array}{ccl}y&=&3x\\z&=&1 \end{array}\right.$.}

      \good {$\Delta$ est contenu dans le plan $z=-2$ ou $z=1$.}

\end{answers}
\begin{explanations}
$\Delta$ est parallèle au plan $(xOy)$ et rencontre l'axe  $(Oz)$, une représentation cartésienne de $\Delta$ est donc  de la forme :
$ \left\{\begin{array}{ccl}ax+by&=&0\\z&=&c, \quad a,b,c \in \Rr \end{array}\right.$.

 On peut supposer que $b$ est non nul, sinon, $\Delta$ ne rencontre pas $D_1$ ou $D_2$. Par conséquent, une représentation cartésienne de $\Delta$ est donc  de la forme :
$ \left\{\begin{array}{ccl}ax+y&=&0\\z&=&b, \quad a,b \in \Rr \end{array}\right.$.

 On calcule $a$ et $b$ pour que $\Delta$ rencontre $D_1$ et $D_2$.
\end{explanations}
\end{question}
}
